\documentclass[a4paper,12pt]{article}
\usepackage[slovene]{babel}
\usepackage[utf8]{inputenc}
\usepackage[T1]{fontenc}
\usepackage{lmodern}

\usepackage{amsmath,amsfonts}
\usepackage{amssymb}

\setlength{\parindent}{0mm}

\begin{document}

% Datoteko dopolnite tako, da bo PDF čim bolj podoben datoteki
% matematicni-izrazi-resitev.pdf.
% Kjer piše ??, manjka matematični izraz.
% Matematičnih nalog ne rešujte, samo zapišite matematične izraze.

\section*{Logika in množice}

\begin{enumerate}
\item
Poišči preneksno obliko formule $\exists x: P(x)\land\forall x: \Rightarrow \forall x : R(x)  $

\item
Definiramo množici $A = \left[2, 5\right]  $ in $B = \{0,1,2,3,4,dots \}$ .
V ravnino nariši:
\begin{enumerate}
\item $(a)  A\cap B\times \emptyset $
\item $(b)  A\cup $
\end{enumerate}

\item
Dokaži:
\begin{itemize}
\item $(A\Rightarrow B) \sim (\neg B \Rightarrow \neg A) $
\item $\neg(A \cup B) \sim \neg A \land \neg B$
\end{itemize}
\end{enumerate}

\section*{Analiza}

\begin{enumerate}
\item
Pokaži, da je funkcija $x \mapsto \sqrt{x}$ enakomerno zvezna na $\left[0,\infty\right] $.

\item
Katero krivuljo določa naslednji parametrični zapis?
% Pomagajte si z dokumentacijo o razmikih v matematičnem načinu:
% https://www.overleaf.com/learn/latex/Spacing_in_math_mode
\[
x(t) = a \cos t,\quad % tu manjka ukaz za presledek
y(t) = b \sin t,\quad % tu manjka ukaz za presledek
t \in [0, 2 \pi]
\]

\item
Pokaži, da ima $f(x)=3x + \sin(2x)$ inverzno funkcijo in izračunaj $(f^-1)(3\pi)$.

\item
Izračunaj integral
% V rešitvi je integral zapisan v vrstičnem načinu, a v prikaznem slogu.
% To dosežete tako, da v matematičnem načinu najprej uporabite ukaz displaystyle.
% Pred dx je presledek: pravi ukaz je \,
$\displaystyle \int_{}^{} \frac{2+\sqrt{x+1}}{(x+1)^2-\sqrt{x+1}} \,dx $
% \frac{2+\sqrt{x+1}}{(x+1)^2-\sqrt{x+1}}

\item
Naj bo \(g\) zvezna funkcija. Ali posplošeni integral
$\int_{0}^{1} \frac{g(x)}{x^2} \,dx $
konvergira ali divergira? Utemelji.

\item
S Taylorjevo formulo pokaži, da v okolici točke \(0\) velja
$\sin x = x- \frac{x^3}{3!} + \mathcal{O} (x^5)$
\end{enumerate}

\section*{Kompleksna števila}

\begin{enumerate}
\item
Naj bo \(z\) kompleksno število, \(z \ne 1\) in $|z|=1 $.
Dokaži, da je število $ i \, \frac{z+1}{z-1} $ realno.

\item
Poenostavi izraz:
% V rešitvi so notranji ulomki zapisani z ukazom dfrac.
$\frac{\dfrac{3+\imath}{2-2\imath}+ \dfrac{7\imath}{1-\imath}}{1+ \dfrac{\imath-1}{4}-\dfrac{5}{2-3\imath}}$
\end{enumerate}

\section*{Zaporedja, vrste in limite}

\begin{enumerate}
\item
Z indukcijo dokaži, da za vsako naravno število \(n\) velja
$\sum_{\imath = 1}^{n}\imath^2 = \frac{n(n+1)(2n + 1)}{6}  $

\item
Naj bo $\sum$ absolutno konvergentna vrsta in \(a_n \ne -1\).
Dokaži, da je tudi vrsta \(\sum_{n=1}^\infty \frac{a_n}{1+a_n}\)
absolutno konvergentna.

\item
Izračunaj limito
$ \lim{\scriptsize x\rightarrow\infty}(\sin\sqrt{x+1}-\sin\sqrt{x})$

\item
Za dani zaporedji preveri, ali sta konvergentni.
% Pomagajte si s spodnjima delno pripravljenima matematičnima izrazoma:
% a_n = \sqrt{2+\sqrt{2+\dots+\sqrt{2}}} \qquad
% b_n = \sin(\sin(\dots(\sin 1)\dots))
% Zaviti oklepaj pod izrazom naredi ukaz underbrace: \underbrace{izraz}_{napis}.
% Besedilo v napisu napišite z ukazom text, presledek pred njim pa z znakom ~.

\qquad $a_n =\underbrace{ \sqrt{2+\sqrt{2+\dots+\sqrt{2}}}}_{n korakov}   \qquad$
$b_n = \underbrace{\sin(\sin(\dots(\sin 1)\dots))}_{n sinusov}$

\end{enumerate}

\section*{Algebra}

\begin{enumerate}
\item
Vektorja $\vec{c}= \vec{a}+ 2\vec{b}$
sta pravokotna in imata dolžino 1. Določi kot med vektorjema \(\vec{a}\) in \(\vec{b}\).

\item
Izračunaj
$$
\begin{pmatrix}
 1 & 2 & 3 & 4 & 5 & 6\\
 4 & 5 & 2 & 6 & 3 & 1
\end{pmatrix}^{-2000}
$$
% 1 2 3 4 5 6
% 4 5 2 6 3 1
\end{enumerate}

\section*{Velika determinanta}

Izračunaj naslednjo determinanto \(2n \times 2n\), ki ima na neoznačenih mestih ničle.
% Prazna mesta v determinanti dobite tako, da med znakoma & ne napišete ničesar.
% Pike naredijo ukazi cdots (vodoravne), vdots (navpične) in ddots (poševne).
$$
\begin{vmatrix}
 1 & & & & 1\\
 & 2 & & & 2\\
&&\ddots & & \vdots\\
&&& n-1  & n-1\\
 1 & 2 &\dots& n-1 & n & n+1&  n+2 & \cdots&2n\\
&&&& n+1 & n+1\\
&&&& n+2 && n+2\\
&&&&\vdots&&&\ddots\\
 &&&&2n &&&& 2n
\end{vmatrix}
$$
\section*{Grupe}

Naj bo
% Definiciji množic G in H sta v rešitvi poravnani pri enačaju, z okoljem align*.
% Zavita oklepaja napišete kot \{ in \}, presledek za podpičjem pa z ukazom \;
$G - \{ z \in \mathbb{C} ; z-2^k(\cos(m\pi\sqrt{2}+\imath\mathbb{C} \sin(m\pi\sqrt{2}))),k,m\in\mathbb{Z})$
$H-\{(x,y)\in \mathbb{R}^2; x,y\in \mathbb{Z} \}$
\begin{enumerate}
\item
Pokaži, da je \(G\) podgrupa v grupi $(\mathbb{C} \backslash  \{0\},\cdot)$
neničelnih kompleksnih števil za običajno množenje.

\item
Pokaži, da je \(H\) podgrupa v aditivni grupi $(\mathbb{R} ^2, +)$
ravninskih vektorjev za običajno seštevanje po komponentah.

\item
Pokaži, da je preslikava \(f:H\to G\), podana s pravilom
$(x,y)\mapsto 2^x(\cos y\pi\sqrt)+(\imath\sin(y\pi\sqrt{2}))$
izomorfizem grup \(G\) in \(H\).
\end{enumerate}

\end{document}
